% ============================================================================
\section{Recomendaciones o sugerencias}

Con base en los resultados empíricos y metodológicos obtenidos, se formulan cinco recomendaciones concretas:
\begin{enumerate}
	\item \textbf{Preservación de archivos históricos con fecha exacta:} conservar bases abiertas de publicaciones sucesivas (\emph{vintages}) del PBI peruano registrando el día calendario exacto de difusión de cada Nota Semanal, para permitir evaluaciones en tiempo real estricto.
	\item \textbf{Documentación transparente de cambios metodológicos:} acompañar las series históricas con metadatos claros que separen las revisiones estadísticas ordinarias de los empalmes por cambio de año base o mejoras de cobertura sectorial.
	\item \textbf{Uso de modelos acotados ante perturbaciones atípicas:} evaluar cualquier sistema de ajuste de revisiones diferenciando épocas estables de episodios extremos, y preferir modelos no lineales acotados (como características aleatorias o árboles potenciados) frente a regresiones lineales irrestrictas cuando existan variaciones bruscas como las de 2020--2021.
	\item \textbf{Cautela institucional frente a modelos cuánticos emergentes:} no incorporar arquitecturas de reservorio cuántico en procesos institucionales de seguimiento económico mientras alternativas clásicas más simples y económicas ofrezcan precisión equivalente o superior fuera de muestra.
	\item \textbf{Transparencia y reproducibilidad:} mantener la práctica de separar un tramo final de prueba intacto y conservar códigos, semillas, parámetros y registros completos de verificación en todo análisis macroeconómico aplicado.
\end{enumerate}

En coherencia con la evidencia negativa controlada observada frente a controles clásicos comparables, una eventual reconsideración futura de modelos cuánticos solo sería pertinente si se contara con series mucho más largas y se demostrara una ganancia estable, estadísticamente significativa y superior en costo y precisión a las alternativas clásicas existentes.
